\tocchapter{Conclusions}

\par This article analysed \nRecords{} documented security weaknesses in the Model Context Protocol ecosystem to find out whether its widely discussed problems come from artificial intelligence or from the speed at which the ecosystem was built. Most of them belong to the second category.

\par For \textbf{RQ1}, \pClassic\% of the records are classic weaknesses that do not need a language model to be exploited, \pNative\% are AI-native and the remaining \pHybrid\% are hybrid. For \textbf{RQ2}, the two main groups have different origins. Classic weaknesses come mostly from implementation errors, which account for half of the dataset, while AI-native weaknesses come from the design of the protocol and of the integrations built on top of it. For \textbf{RQ3}, the clearest signs of rushed work are insecure default settings, present in more than a quarter of the classic records, errors in first-party reference code that other developers reuse, security features added to the specification after it was already in use, and the installation of servers without any review. None of the AI-native records received a \acrshort{glo:cve} identifier, so tools and processes that rely on vulnerability databases see only the classic part of the risk. The rest has to be handled through threat modelling and system design.

\par For software engineers, the practical conclusion is that most of these failures could have been avoided with standard practices: no shell strings built from user input, proper path checks, authentication enabled by default, filtering of outgoing requests and review of dependencies. The Moldovan e-government breach shows that a related mistake, trusting the calling application instead of checking the user, causes serious damage even without AI. Institutions that plan to connect AI assistants to public registries should first change this trust model.

\par \textbf{Future work.} A second, independent coder should classify these records so that agreement can be measured. The dataset should also be extended beyond \lastDisclosure{} and compared with the \acrshort{glo:cve} dataset of He et al. Finally, the proposed controls could be tested on a sample of open-source servers.
